\chapter{Фаза детерминанта и вещественная масса: тест критерия Ф.2}

После выбора anomaly inflow как вспомогательной invertible/topological
теории остаётся один решающий вопрос: может ли eta-фаза фермионного
детерминанта породить вещественный аддитивный член \(-\alpha/3\) в
относительном операторе заряженных лептонов?

\section{Полный product Dirac operator}

Для product geometry квадрат оператора имеет вид
\begin{equation}
 D_K^2=D_{\mathbb{RP}^3}^2+D_{S^1}^2,
\end{equation}
поэтому каждому трёхмерному собственному значению \(\lambda_j\) и окружной
моде \(n+\beta\) соответствует положительная энергия
\begin{equation}
 E_{j,n}^2=\lambda_j^2+\frac{(n+\beta)^2}{R_1^2}.
\end{equation}
Физический charged lepton является vectorlike Dirac field. Его полный
Clifford-блок содержит пару \(\pm iE_{j,n}\), и при вещественной массе
\begin{equation}
 (m+iE_{j,n})(m-iE_{j,n})=m^2+E_{j,n}^2>0.
\end{equation}
Следовательно,
\begin{equation}
 \log(m+iE)+\log(m-iE)=\log(m^2+E^2),
\end{equation}
а mass response равен
\begin{equation}
 \frac{\partial}{\partial m}\log\det(D_K+m)
 =\sum_{j,n}\frac{2m}{m^2+E_{j,n}^2}.
\end{equation}
Обе формулы вещественны blockwise. Спектральная асимметрия нечётномерного
фактора не создаёт eta-члена в производной вещественного vectorlike
effective action.

\section{Численный asymmetric-spectrum control}

Машинный тест намеренно использует несимметричный набор положительных и
отрицательных трёхмерных eigenvalues, после чего строит полный
четырёхмерный Clifford pair для periodic и antiperiodic окружностей. Для
каждого блока, полной суммы и antiperiodic-minus-periodic difference мнимая
часть равна нулю с машинной точностью. Производная по массе также полностью
вещественна. Тем самым cancellation не опирается на случайную симметрию
выбранного cutoff spectrum.

\section{Почему \texorpdfstring{\(\zeta(-1)\)}{zeta(-1)} не является этим
детерминантом}

Для одного massive product-level конечное circle determinant ratio задаётся
\begin{equation}
 \mathcal I_\rho(\beta)=
 \log\frac{\cosh(2\pi\rho)-\cos(2\pi\beta)}
 {\cosh(2\pi\rho)-1}.
\end{equation}
При \(\rho=1\) и \(\beta=1/2\) получаем
\begin{equation}
 \mathcal I_1(1/2)=0.0074697796\ldots,
\end{equation}
что точно воспроизводит прежний projected-KK audit, но не равно
\begin{equation}
 |\zeta_R(-1)|=\frac1{12}=0.0833333\ldots.
\end{equation}
Первое число является mass-dependent Euclidean log-determinant; второе ---
Casimir-регуляризацией линейной суммы уровней. Их нельзя складывать как
части одного determinant без нового functional bridge.

\section{Статус трёхмерной eta-фазы}

В текущем spin-аудите для двух структур \(\mathbb{RP}^3\) зафиксировано
\begin{equation}
 \eta_D=\pm\frac14.
\end{equation}
Значение \(-1/2\) отсутствует в замороженном operator menu и потребовало бы
явного twisted или doubled Dirac operator. Но даже после такого изменения
eta-фаза относится к chiral determinant line. В vectorlike charged-lepton
determinant две хиральности образуют положительный модуль и фазовый вклад
сокращается.

\begin{theorem}[Минимальный phase-to-mass no-go]
Для вещественной положительной Dirac-массы и vectorlike charged-lepton
operator на \(\mathbb{RP}^3\times S^1\) eta-инвариант не даёт аддитивного
вклада в \(\partial_m\Re\log\det(D_K+m)\). Поэтому разложение
\begin{equation}
 \frac13=\frac14+\frac1{12}
\end{equation}
не выводит вещественный mass shift \(-\alpha/3\) в минимальной Gaussian
реализации.
\end{theorem}

\begin{proof}
Каждый product-level входит парой \(m\pm iE\). Их произведение положительно,
мнимая часть логарифма сокращается, а mass derivative зависит только от
\(m^2+E^2\). Отдельно, окружной determinant является функцией
\(\mathcal I_\rho(\beta)\), а не константой \(|\zeta_R(-1)|\). Следовательно,
ни eta-фаза, ни Casimir-остаток не поставляют требуемый вещественный
коэффициент в объявленном operator class.
\end{proof}

\begin{protocol}[Условие переоткрытия Ф.2]
Phase-to-mass ветвь может быть переоткрыта только если из parent action
следуют:
\begin{enumerate}
 \item complex chiral Yukawa mass с физически неустранимой фазой;
 \item CP-odd background или сумма topological sectors, превращающая
 фазовую интерференцию в вещественный potential;
 \item заранее фиксированная связь этого potential с lepton mass modulus;
 \item линейный коэффициент \(\alpha/3\) без использования \(m_\tau\);
 \item ablation eta/topological sector уничтожает именно вещественный mass
 shift, а не только фазу partition function.
\end{enumerate}
\end{protocol}

Текущий вердикт критерия Ф.2 отрицателен. 4D topological interpretation
остаётся допустимой для global anomaly и determinant-phase bookkeeping, но
не закрывает \(\tau/\mu\)-формулу. Единственным существующим вещественным
кандидатом остаётся прежняя QED winding self-energy, для которой unit
normalization даёт коэффициент \(0.0616969\ldots\), а требуемый projection
weight \(5.4027533\ldots\) не выведен.

Машинный протокол сохранён в
\texttt{s2t\_eta\_phase\_mass\_gate\_results.json}.